%! Author = lili
%! Date = 4/27/26


\section{Vorlesung 4}
\begin{bsp}
    \textbf{n-faches Werfen einer fairen Münze} \\
    \[
        \glssymbol{eraum} _n = \kl{0,1}^n, \quad \glssymbol{wfk}_n (\kl{\glssymbol{elementarereignis}}) = (\frac{1}{2})^n \faoio \txt{(Gleichverteilung)}
    \]
    $\glssymbol{elementarereignis}_i = 1$ bedeute ``Kopf'' \\
    Falls es uns interessiert, wie oft ``Kopf'' fällt:
    \[
        S_n(\glssymbol{elementarereignis}) \coloneqq  \sum^n_{i=1}     \glssymbol{elementarereignis}_i \txt{ist Anzahl ''Erfolge'' (Kopf)}
    \]
    \[
        S_n :\glssymbol{eraum} _n \rightarrow \kl{0,1,..,n} \txt{Abbildung!}
    \]
    Falls der relative Anteil der Erfolge relevant ist:
    \[
        T_n: \glssymbol{eraum} _n \rightarrow [0,1]
    \]
    \[
        T_n (\glssymbol{elementarereignis}) \coloneqq  \frac{1}{n} S_n (\glssymbol{elementarereignis}) = \frac{1}{n} \sum^n_{i=n} \glssymbol{elementarereignis}_i \faoio
    \]
\end{bsp}

\begin{defi}
    \label{def:zv}
    \textbf{\gls{zv}} \\
    Für einen \gls{wraum} $\wraum$ und eine Menge $E \neq \emptyset$ nennen wir eine messbare Abbildung
    \[
        X: \glssymbol{eraum}  \rightarrow E
    \]
    eine $E$-wertige \gls{zv}.
\end{defi}
\komm{\defirefman{zv}: \gls{zv} ist keine Variable: Es ist eine Zuordnungsvorschrift der Stochastik, welche jedem möglichen Ergebnis eines Zufallsexperiments eine Größe zuordnet. Also das Ergebnis eines Zufallsexperiments.}

\begin{defi}
    \label{def:verteilung-zv}
    \textbf{Verteilung einer \gls{zv}}\\
    Sei $X$ eine $E$-wertige \gls{zv} auf $\wraum$
    \[
        \pe_X(B) \coloneqq  \pe(\kl{\oio : X(\glssymbol{elementarereignis}) \in B}) \quad \forall \txt{messbaren } B \subset E = \oub{\pe(X \in B)}{Kurzschreibweise}
    \]
    Das \gls{wmaß} $\pe_X$ auf $E$ heißt Verteilung von $X$.
\end{defi}
\komm{\defirefman{verteilung-zv}: Die Verteilung von X beschreibt, wie die Werte von X verteilt sind, also
\textbf{wie wahrscheinlich verschiedene Ergebnisse von X sind}.}
\begin{lem}
    Ist $\glssymbol{eraum} $ abzählbar und $\pe$ durch die \gls{wfk} $\glssymbol{wfk}$ gegeben, so ist $\pe_X$ ein \gls{wmaß}.

    \paragraph{Beweis.}
    \[
        \glssymbol{wfk}_X (e) = \pe_X(\kl{e}) = \txc{\pe(X \in \kl{e}) = } \pe(X=e) = \sum_{\substack{\oio \\ X(\glssymbol{elementarereignis}) = e}} \glssymbol{wfk}(e) \geq 0
    \]
    z.z.\ Normierung
    \[
        \sum_{e \in E} \glssymbol{wfk}_X (e) = \sum_{e \in E} \sum_{\substack{\oio :\\ X(\glssymbol{elementarereignis}) = e}} \glssymbol{wfk} (\glssymbol{elementarereignis}) =  \sum_{\substack{\oio \\ e \in E :\\ X(\glssymbol{elementarereignis}) = e}} \glssymbol{wfk}  (\glssymbol{elementarereignis}) = \sum_{\substack{\oio :\\ X(\glssymbol{elementarereignis}) = e}} \glssymbol{wfk} (\glssymbol{elementarereignis}) = \sum_{\oio} \glssymbol{wfk} (\glssymbol{elementarereignis}) = 1
    \]
\end{lem}


\section{Notationen}
\[
    \pe(X \in B) = \pe(\kl{\oio: X(\glssymbol{elementarereignis}) \in B})
\]
Falls $B$ ein Intervall ist:
\[
    \pe(X \geq b), \quad \pe(a < X \leq b) , \dots
\]
\begin{defi} \label{def:wert-bild-e}
    Geg.:
    \[
        X: \glssymbol{eraum}  \rightarrow E \txt{\gls{zv} auf } \wraum
    \]
    Dann heißt $E$ der Wertebereich und $X(\glssymbol{eraum} )\coloneqq  \kl{X(\glssymbol{elementarereignis}): \oio}$ heißt Bildbereich.
\end{defi}
\begin{bsp} \defirefman{wert-bild-e}
    \txo{
        \[
            S_n : \glssymbol{eraum} _n \rightarrow \kl{0,\dots,n}
        \]
        \[
            S_n : \glssymbol{eraum} _n \rightarrow \mathbb{N}_0
        \]
    }
\end{bsp}
\begin{bem} \defirefman{wert-bild-e} \\
    $X(\glssymbol{eraum} ) \subset E$ gilt immer.
\end{bem}
\komm{Denn $X(\glssymbol{eraum} ) $ ist ein $\oio$ und jedes \glssymbol{elementarereignis} ist ein Element aus $E$.}
\begin{bem} \defirefman{wert-bild-e} \\
    Gegeben ist ein \gls{wmaß} auf $\tilde {\glssymbol{eraum}}$, so existiert eine \gls{zv}, deren Verteilung genau dieses \gls{wmaß} ist. \\
    \noindent Sei $(\tilde {\glssymbol{eraum}} , \tilde\pe)$ ein beliebiger \gls{wraum}. \\
    $X: \glssymbol{eraum}  \rightarrow \tilde{\glssymbol{eraum}} $ definiert durch:
    \begin{align*}
    {\glssymbol{eraum}}
        & \coloneqq  \tilde{\glssymbol{eraum}}  \\
        \pe&\coloneqq  \tilde\pe\\
        X(\glssymbol{elementarereignis}) & = \glssymbol{elementarereignis} \txt{Identität} \faoio
    \end{align*}
    \[
        \pagebreak(X \in B) = \pe(\oub{\kl{\oio: X(\glssymbol{elementarereignis}) \in B}}{=B}) = \pe(B) = \tilde \pe(B)
    \]
    $\Rightarrow X$ hat Verteilung $\tilde\pe$
\end{bem}
\komm{Zu jedem Wahrscheinlichkeitsmaß auf einem Messraum existiert eine Zufallsvariable, die genau dieses Maß als Verteilung hat. \\
aka. Für jede vorgegebene Wahrscheinlichkeitsverteilung kann man ein Zufallsexperiment konstruieren, das diese Verteilung als Ergebnis hat.}


\begin{defi}
    \textbf{\dichte einer \gls{zv}} \\
    Gegeben \dichte \glssymbol{dichte} \\
\end{defi}

\paragraph{Hier fehlt was}
\newpage
